\section{Answer to the Research Question}
\label{sec:conclusion_summary}

This thesis asked whether conventional institutional ownership measures and the topology of institutional ownership improve the out-of-sample prediction of stock-level downside risk beyond market characteristics, whether graph-based models exploit that structure more effectively than comparable tabular models, and whether the resulting predictions have economic value. The answer is qualified. Conventional ownership measures contain limited incremental information, but neither the engineered network features nor direct graph learning provide a consistent advantage sufficient to justify their additional complexity. The strongest result is instead that the selected model's predictions provide an informative out-of-sample ranking of subsequent downside risk.

The empirical design separates information, model flexibility, and evaluation timing. Quarterly Form 13F holdings are aligned with CRSP market data using a common information cutoff, and the primary outcome is the positive magnitude of maximum drawdown over the following 63 trading days. Market, ownership, and engineered network features form nested information sets evaluated using Ridge and XGBoost, while three graph neural networks learn directly from the manager--stock graph. Feature selection, preprocessing, tuning, and model comparison use chronological training and validation samples separated from the 2024~Q1--2025~Q4 test period by purged quarters. The test sample is examined only after the model-selection decision is fixed.

The feature comparisons provide limited support for Hypothesis~H1: ownership variables produce small validation improvements, but not enough to displace the market-only specification under the stated parsimony rule. Engineered network features add no consistent improvement, so Hypothesis~H2 is not supported. Market-only XGBoost modestly outperforms market-only Ridge on validation and test data, supporting Hypothesis~H3. At the validation stage, Temporal GraphSAGE has the lowest MAE, whereas XGBoost has better RMSE and $R^2$ and essentially identical ranking performance. Market-only XGBoost is therefore selected as the performance--parsimony choice before the test sample is examined. The held-out comparison preserves this pattern: Temporal GraphSAGE has the lowest MAE, while XGBoost has better RMSE and $R^2$ and nearly identical ranking performance. The graph evidence is competitive but not consistently superior, giving mixed support to Hypothesis~H4.

On the held-out sample, the selected XGBoost model achieves an MAE of 0.0912, an $R^2$ of 0.497, and a mean quarterly Spearman correlation of 0.738. Mean realized maximum drawdown increases monotonically across Fragility Score deciles in every test quarter, and the average Decile~10--Decile~1 spread is 41.6 percentage points. Positive rank correlations and decile spreads across the alternative targets show that the ranking extends beyond the primary 63-day drawdown definition. The error analysis nevertheless shows that forecasts understate some of the most severe realized losses. The score is therefore better supported as a relative measure of downside exposure than as a precisely calibrated forecast of extreme drawdown magnitude.

The portfolio evidence is more conditional. Excluding the most fragile decile improves the downside profile of the matched equal-weighted universe, and the screened-minus-unscreened return difference has a positive FF5-plus-momentum alpha. The screened portfolio itself has approximately zero alpha, does not outperform the external CRSP total-market benchmark, and provides almost no improvement under value weighting. The exploratory long--short portfolios further separate low- and high-score stocks, but the short evaluation period, substantial turnover, large drawdowns, and incomplete treatment of short-side costs do not support a tradable-factor interpretation. Hypothesis~H5 consequently receives qualified support. The most defensible application of the Fragility Score is risk monitoring and portfolio screening, particularly within a specified investment universe, rather than return forecasting or stand-alone trading.
